\clearpage
\noindent\makebox[10mm][l]{42}\parbox[t]{\dimexpr\textwidth-10mm\relax}{\centering\scriptsize DETERMINATION OF THE DIFFERENCE IN LATITUDE BETWEEN GÖTTINGEN AND ALTONA.}\par
1820, reckoned from the lower culmination of 13 May:
\begin{center}May 13.\quad Upper culm.\ $-0''{,}10$;\end{center}
1824, reckoned from the upper culmination of 20 April:
\[
\begin{array}{rl|r}
\text{Apr.}&21.\ \text{L. C.}&-0''{,}13\\
&21.\ \text{U. C.}&-0{,}26\\
&25.\ \text{L. C.}&-1{,}29\\
&27.\ \text{U. C.}&-2{,}04\\
&28.\ \text{U. C.}&-2{,}32\\
&29.\ \text{L. C.}&-2{,}45\\
\text{May}&1.\ \text{L. C.}&-2{,}93\\
&1.\ \text{U. C.}&-3{,}03\\
&2.\ \text{L. C.}&-3{,}14\\
&8.\ \text{U. C.}&-4{,}64\\
&9.\ \text{L. C.}&-4{,}77
\end{array}
\]
\begin{center}14.\end{center}
Let $f$ and $g$ denote the flexure of the telescope with its optical axis horizontal and vertical, respectively. By flexure is meant the change in position of the optical axis projected onto the plane of the graduated circle, relative to the graduations, owing to the effect of gravity on all the connected parts of the instrument. If we assume that this flexure is proportional to the force of gravity (an assumption that seems safe given the extremely small magnitude of the whole effect), then for an inclination $z$ of the optical axis the flexure is expressed by $f\sin z+g\cos z$. Thus, if the collimation error is $=e$ and the zenith distance as read is $=z$, the true zenith distance
\[
=z-e+f\sin(z-e)+g\cos(z-e)
\]
is obtained. If the telescope were perfectly symmetrical, $g$ would vanish entirely. But since no material object made by human hands is absolutely perfect, and since the balancing weights themselves already disturb perfect symmetry to some extent, it seems quite reasonable to admit that one or a couple of tenths of a second might be the

\clearpage
\noindent\mbox{}\hfill III.\quad RESULTS.\hfill 43\par
\noindent value of $g$. Once the calculation is carried to individual tenths or even hundredths of a second, it would be inconsistent to omit the second part of the flexure where it can be taken into account.
\begin{center}15.\end{center}
The complement to $90^\circ$ of half the difference between the zenith distances measured directly and by reflection gives the zenith distance freed from the collimation error and from the first part of the flexure. It thus contains only the second part of the flexure, with opposite signs [depending on] whether the circle is to the east or to the west. This zenith distance evidently refers to the vertical at the point where the optical axis meets the water vessel. That vertical differs imperceptibly between the two culminations of Polaris and is $0''{,}05$ farther north than the axis of the circle. This combination also better suits our purpose in that it avoids the assumption that the collimation error remained constant throughout the observations of 1824. If the number of direct observations is $=\alpha$ and the number of observations by reflection is $=\beta$, the weight of that determination is $=\dfrac{4\alpha\beta}{\alpha+\beta}$, insofar as the observational errors are regarded as purely random and independent of one another.
\begin{center}16.\end{center}
If we now denote by
\begin{quote}
\begin{tabular}{@{}l@{\quad}p{.84\textwidth}@{}}
$\varphi$&the altitude of the pole at the site of the water vessel\\
$\delta$&the declination of Polaris at the lower culmination of 13 May 1820\\
$\delta'$&the declination at the upper culmination of 20 April 1824,
\end{tabular}
\end{quote}
then the observations give us the following determinations:
\begin{center}
\begin{tabular}{lllcll}
\multicolumn{6}{l}{for $\delta+\varphi-0{,}765g$}\\
\quad1820&May 13&$\ldots\,\ldotp$&$139^\circ52'38''{,}40$&Weight&$6{,}86$\\
\multicolumn{6}{l}{for $\delta-\varphi+0{,}800g$}\\
\quad1820&May 13&$\ldots\,\ldotp$&$36^\circ49'\phantom0 1''{,}50$&Weight&$2{,}00$
\end{tabular}
\end{center}
\noindent\hfill $6^*$\par

\clearpage
\noindent\makebox[10mm][l]{44}\parbox[t]{\dimexpr\textwidth-10mm\relax}{\centering\scriptsize DETERMINATION OF THE DIFFERENCE IN LATITUDE BETWEEN GÖTTINGEN AND ALTONA.}\par
\begin{center}
\begin{tabular}{rrrc rr r c r}
\multicolumn{9}{l}{for $\delta'+\varphi-0{,}765g$}\\
1824&Apr.&21&$\ldots\,\ldotp$&$139^\circ$&$54'$&$5''{,}61$&Weight&$6{,}86$\\
&&25&$\ldots\,\ldotp$&139&54&$5{,}76$&»&$6{,}86$\\
&&29&$\ldots\,\ldotp$&139&54&$6{,}36$&»&$6{,}86$\\
&May&1&$\ldots\,\ldotp$&139&54&$5{,}91$&»&$6{,}86$\\
\multicolumn{9}{l}{for $\delta'-\varphi+0{,}800g$}\\
1824&Apr.&20&$\ldots\,\ldotp$&$36^\circ$&$50'$&$31''{,}15$&Weight&$2{,}67$\\
&&21&$\ldots\,\ldotp$&36&50&$30{,}29$&»&$6{,}86$\\
&&27&$\ldots\,\ldotp$&36&50&$30{,}65$&»&$6{,}86$\\
&&28&$\ldots\,\ldotp$&36&50&$30{,}73$&»&$6{,}86$\\
&May&1&$\ldots\,\ldotp$&36&50&$30{,}25$&»&$6{,}86$\\
\multicolumn{9}{l}{for $\delta'+\varphi+0{,}765g$}\\
1824&May&2&$\ldots\,\ldotp$&$139^\circ$&$54'$&$6''{,}71$&Weight&$6{,}86$\\
&&9&$\ldots\,\ldotp$&139&54&$6{,}15$&»&$6{,}86$\\
\multicolumn{9}{l}{for $\delta'-\varphi-0{,}800g$}\\
1824&May&8&$\ldots\,\ldotp$&$36^\circ$&$50'$&$30''{,}48$&Weight&$6{,}86$.
\end{tabular}
\end{center}
We thus obtain the following six equations for determining the four unknown quantities $\delta$, $\delta'$, $\varphi$, $g$
\[
\begin{array}{rclrrrcr}
\delta+\varphi-0{,}765g&=&139^\circ&52'&38''{,}40&&\text{Weight}&6{,}86\\
\delta-\varphi+0{,}800g&=&36&49&1{,}50&&\text{»}&2{,}00\\
\delta'+\varphi-0{,}765g&=&139&54&5{,}91&&\text{»}&27{,}43\\
\delta'-\varphi+0{,}800g&=&36&50&30{,}54&&\text{»}&30{,}10\\
\delta'+\varphi+0{,}765g&=&139&54&6{,}43&&\text{»}&13{,}71\\
\delta'-\varphi-0{,}800g&=&36&50&30{,}48&&\text{»}&6{,}86.
\end{array}
\]
from which the method of least squares\textsuperscript{*)} gives the following values:
\[
\begin{array}{rcrrr}
\delta&=&88^\circ&20'&50''{,}33\\
\delta'&=&88&22&18{,}28\\
\varphi&=&51&31&47{,}90\\
g&=&\multicolumn{3}{r}{+\ 0{,}17.}
\end{array}
\]
The weight of the determination of $\varphi$ is then $=60{,}8$.
\par\noindent\rule{.25\textwidth}{.4pt}\par
{\footnotesize *) Somewhat more conveniently here by the procedure in Supplem.\ Theor.\ Comb.\ Observ.\par}

\clearpage
\noindent\mbox{}\hfill III.\quad RESULTS.\hfill 45\par
To obtain some measure of the accuracy of the observations, we substitute these values into the fourteen equations from which the preceding six were derived. The following errors then remain:
\begin{center}
\begin{tabular}{r|c}
Error&Weight of equation\\\hline
$-0''{,}30$&$6{,}86$\\
$+1{,}07$&$2{,}00$\\
$+0{,}44$&$6{,}86$\\
$+0{,}29$&$6{,}86$\\
$-0{,}31$&$6{,}86$\\
$+0{,}14$&$6{,}86$\\
$-0{,}63$&$2{,}67$\\
$+0{,}23$&$6{,}86$\\
$-0{,}13$&$6{,}86$\\
$-0{,}21$&$6{,}86$\\
$+0{,}27$&$6{,}86$\\
$-0{,}40$&$6{,}86$\\
$+0{,}16$&$6{,}86$\\
$-0{,}24$&$6{,}86$
\end{tabular}
\end{center}
The sum of the products of the squares of these errors and their weights is $=9{,}6154$. Hence an approximate value for the mean error of an observation
\[
=\sqrt{\frac{9{,}6154}{10}}=0''{,}981.
\]
The mean error to be expected in the final result for the altitude of the pole, insofar as it arises from causes acting irregularly, is therefore
\[
=\frac{0''{,}981}{\sqrt{60{,}8}}=0''{,}126.
\]
The uncertainty of the result must, however, be somewhat greater, since the assumption that all observational errors occur without pattern and are independent of one another is not quite correct. For observa-

\clearpage
\noindent\makebox[10mm][l]{46}\parbox[t]{\dimexpr\textwidth-10mm\relax}{\centering\scriptsize DETERMINATION OF THE DIFFERENCE IN LATITUDE BETWEEN GÖTTINGEN AND ALTONA.}\par
\noindent tions of the same kind at one culmination, and for culminations of the same kind on several days, the underlying reading is practically the same. Although readings taken with verniers almost always bring different graduation marks into play, so that in our procedure their irregular graduation errors are included in the mean error of an observation, $0''{,}981$, it is natural that different average errors should prevail in different parts of the limb. In any event they are very small. In 1826 I used four excellent microscopes by \textsc{Repsold} to examine 30 graduation marks at intervals of 12 degrees with the utmost care, setting each mark almost 200 times in different combinations. The result is that the mean of the errors of two diametrically opposite graduation marks, $A$ and $A+180^\circ$, insofar as any regularity can still be discerned, is represented as closely as possible by the formula
\[
-1''{,}23\cos(2A-28^\circ28')-0''{,}22\cos(4A-47^\circ56')
\]
while the errors then remaining appear irregular, and the square root of the mean of their squares is $=0''{,}32$. I had intended to extend this examination to twice as many graduation marks. Given the small magnitude of the results, however, the examination hardly seems to warrant the considerable time it would require. The first part of the regular error, $-1''{,}23\cos(2A-28^\circ28')$, of course cancels of its own accord when all four verniers are read, as was always done in the observations above. It does, however, provide an actual correction if the graduations are read at only two opposite points, as I now always do, having gained considerably in the precision of my readings by using two \textsc{Repsold} microscopes in place of the verniers.
\begin{center}17.\end{center}
If one prefers to assume $g=0$, the altitude of the pole comes out $0''{,}07$ smaller, and the weight of this determination is $=84{,}1$. Other observations, to be reported elsewhere, do appear to confirm the above

\clearpage
\noindent\mbox{}\hfill III.\quad RESULTS.\hfill 47\par
\noindent value of $g$ in sign and very nearly also in magnitude, but are still insufficient to decide so delicate a matter.

The coefficient $f$ cannot be determined from the present observations without assuming that the collimation error remained constant during the observations of 1824. If we allow ourselves this assumption, we have 28 equations, and treating them in the appropriate way gives
\[
\begin{array}{rcll}
\varphi&=&51^\circ31'47''{,}89&\text{with weight }60{,}9\\
f&=&\phantom{51^\circ31'}+\ 0{,}76&\\
g&=&\phantom{51^\circ31'}+\ 0{,}23&
\end{array}
\]
Since the collimation error can now be determined every hour with remarkable accuracy, without reversal, by setting the telescope on the nadir point\textsuperscript{*)}, I reserve the matter for further examination.
\begin{center}18.\end{center}
Subject to a correction that further investigations have yet to determine and that can scarcely amount to half a second, I therefore take the altitude of the pole to be
\begin{center}
\begin{tabular}{p{.72\textwidth}r}
in Göttingen&\\
\quad at the site of the water vessel used in observing Polaris&$51^\circ31'47''{,}90$\\
\quad at the site of \textsc{Reichenbach}'s meridian circle&$47{,}85$\\
\quad at the site of the zenith sector&$47{,}92$\\
\multicolumn{2}{p{.94\textwidth}}{\quad\quad (the latter being used as the basis for reducing the declinations of the zenith stars)}\\
in Altona&\\
\quad at the site of the zenith sector&$53^\circ32'44''{,}42$\\
\quad at the site of the meridian circle&$45{,}27$.
\end{tabular}
\end{center}
\par\noindent\rule{.25\textwidth}{.4pt}\par
{\footnotesize *) I have used this invaluable means continually for two years. \textsc{Bohnenberger} first showed that it was practicable.\par}

\clearpage
\noindent\makebox[10mm][l]{48}\parbox[t]{\dimexpr\textwidth-10mm\relax}{\centering\scriptsize DETERMINATION OF THE DIFFERENCE IN LATITUDE BETWEEN GÖTTINGEN AND ALTONA.}\par
\begin{center}19.\end{center}
According to the trigonometrical connection between the observatories of Göttingen and Altona, the latter lies
\[
\begin{array}{rl}
115163{,}725&\text{toises north}\\
7{,}211&\text{toises west}
\end{array}
\]
of the former. These figures refer to the sites of the meridian circles. They are based on the value of the Hamburg-Hohenhorn side of the triangle, $13841{,}815$ toises, which in turn is based on the baseline measured in Holstein in 1820 by Prof.\ \textsc{Schumacher}. Since the comparison of the measuring rods used for this purpose with the standard toise has yet to be \textls[150]{definitively} completed, the above distance will eventually have to be changed in the same proportion as the baseline itself, though in any event this change can only be very small. The mean length of a degree of latitude between the two observatories accordingly comes out
\[
=57127{,}2\text{ toises},
\]
appreciably greater than one would have expected from the mean values of the degrees measured in France and England.
\begin{center}20.\end{center}
The Hanoverian arc measurement thus provides further confirmation of the now indubitable fact that the surface of the Earth is not altogether regular in shape. The anomalies in parts of the French and English arc measurements have already supplied evidence of this irregularity, and the anomalies in the altitudes of the pole at several places in Italy have supplied still stronger evidence. In the Hanoverian arc measurement, besides the anomaly between Göttingen and Altona, there is a considerably greater one at an intermediate triangulation point, the Brocken. If my triangles are computed as lying on the surface of an elliptical spheroid, with the dimensions that \textsc{Walbeck} derived from all the arc measurements made up to that time, and which to the best of our present knowledge most closely fits the actual shape

\clearpage
\noindent\mbox{}\hfill III.\quad RESULTS.\hfill 49\par
\noindent\textls[150]{as a whole} (flattening $\dfrac1{302{,}78}$, the three hundred and sixtieth part of the Earth's meridian $=57009{,}758$ toises), and if the altitude of the pole at Göttingen $=51^\circ31'47''{,}85$ is taken as the starting point, one finds the latitude
\[
\begin{array}{lcrrr}
\text{of the Brocken}&=&51^\circ&48'&1''{,}85\\
\text{of Altona}&=&53&32&50{,}79.
\end{array}
\]
Now whereas the astronomical observations gave the altitude of the pole at Altona as $5''{,}52$ \textls[150]{smaller}, the observations made by Mr.\ von \textsc{Zach} on the Brocken give the altitude of the pole at that point as $10\text{—}11''$ \textls[150]{greater}\textsuperscript{*)}, a difference of which in any case only a small part can be ascribed to the instrument and the declinations used in the calculation. Comparing the difference in latitude between Altona and the Brocken with the curvature corresponding to the spheroid that best fits the Earth as a whole would therefore give a discrepancy of $16''$.

In our view, to speak in such cases only of local deflections of the plumb line, as if these phenomena were merely isolated exceptions, is to look at the matter from the wrong point of view. What we call the surface of the Earth in the geometrical sense is simply the surface that everywhere intersects the direction of gravity at right angles and of which the surface of the oceans forms part. The direction of gravity at each point is determined by the shape of the solid part of the Earth and its uneven density. At the outer crust of the Earth, the only part we know anything about, this shape and density prove highly irregular. The irregularity in density may well extend to a considerable depth beneath the outer crust and lies entirely beyond our calculations, for which almost all the data are lacking. The geometrical surface is the product of the combined effect of these unevenly distributed elements. When unmistakable evidence of irregularity presents itself, the greater wonder would seem to be that the irregularity is not greater still.
\par\noindent\rule{.25\textwidth}{.4pt}\par
{\footnotesize *) Monatl.\ Corresp.\ vol.\ X.\ p.\ 203. At a site about $0''{,}5$ south of the triangulation point, this skilled observer found, from 188 observations of $\alpha$ Aquilae, $51^\circ48'12''{,}12$. From observations of the Sun he found $51^\circ48'11''{,}17$.\par}
\noindent IX.\hfill 7\par

\clearpage
\noindent\makebox[10mm][l]{50}\parbox[t]{\dimexpr\textwidth-10mm\relax}{\centering\scriptsize DETERMINATION OF THE DIFFERENCE IN LATITUDE BETWEEN GÖTTINGEN AND ALTONA.}\par
\noindent If astronomical observations were capable of ten or a hundred times their present accuracy, they would no doubt detect this irregularity everywhere.

As matters stand, nothing prevents us from regarding the Earth as a whole as an elliptical spheroid of revolution, from which the actual (geometrical) surface departs everywhere in undulations that are sometimes stronger, sometimes weaker, sometimes shorter, sometimes longer. If it were possible to cover the whole Earth, as it were, with a single trigonometrical network and thereby compute the relative positions of all points, the ideal spheroid of revolution would be the one on which the computed directions of the verticals would agree as closely as possible with the astronomical observations. Although we will always remain far from this unattainable ideal, there is no doubt that future centuries will be able to advance our mathematical knowledge of the figure of the Earth very considerably. Multiplying the number of arc measurements is really only the beginning, yielding only individual results for a small number of points on isolated lines. How much richer the yield will be if the trigonometrical operations carried out in various countries with carefully chosen means are linked and brought together into a single large system. Perhaps the prospect that all the observatories in Europe will one day be connected trigonometrically is not fanciful, since such connections already exist from Scotland to the Adriatic Sea and from Formentera to Funen, although thus far only some of them have been made public. May this last circumstance receive more attention than it has so far, and may valuable materials that should belong to the scientific world neither be withheld from it nor even be exposed to the risk of being lost!
\begin{center}21.\end{center}
Comparing the stellar declinations obtained from the sector observations with earlier determinations, where such exist, yields a result of some interest. Of our 43 stars,

\clearpage
\noindent\mbox{}\hfill III.\quad RESULTS.\hfill 51\par
\noindent 27 are in \textsc{Piazzi}'s catalogue and 13 in \textsc{Bessel}'s \textsc{Bradley} catalogue. Here follows the comparison of our determinations (1827) with \raisebox{.3ex}{.}those of \textsc{Bradley} (1755) and \textsc{Piazzi} (1800), reduced according to \textsc{Bessel}'s new determination of precession. Positive signs indicate a more northerly position according to our determination.
\begin{center}
{\small\setlength{\tabcolsep}{4pt}
\begin{tabular}{r|l|r|r||r|l|r|r}
\multicolumn{2}{c|}{Designation}&\textsc{Bradley}&\textsc{Piazzi}&\multicolumn{2}{c|}{Designation}&\textsc{Bradley}&\textsc{Piazzi}\\\cline{1-4}\cline{5-8}
1&24 Canum&$+0''{,}2$&$+1''{,}1$&17&P.\ 15.39&—&$-0''{,}7$\\
2&83 Ursae&$-1{,}5$&$-2{,}0$&25&$\theta$ Draconis&$+23''{,}5$&$+8{,}6$\\
3&$\eta$ Ursae&$-2{,}4$&$-2{,}3$&27&P.\ 16.33&—&$-2{,}6$\\
4&86 Ursae&$-5{,}4$&$-0{,}8$&28&P.\ 16.56&—&$-2{,}6$\\
6&P.\ 13.289&—&$-1{,}3$&32&16 Draconis&$+1{,}2$&$-3{,}7$\\
7&13 Bootis&$+1{,}4$&$+2{,}4$&36&P.\ 16.253&—&$-2{,}1$\\
8&$\varkappa$ Bootis sq.&$-2{,}9$&$-2{,}2$&37&P.\ 16.291&—&$+10{,}3$\\
9&P.\ 14.56&—&$-0{,}4$&38&P.\ 16.310&—&$-4{,}9$\\
10&$\theta$ Bootis&$-30{,}6$&$-10{,}8$&39&P.\ 17.20&—&$-3{,}0$\\
11&P.\ 14.131&—&$+7{,}4$&40&P.\ 17.38&—&$-2{,}9$\\
12&P.\ 14.164&—&$+0{,}1$&41&74 Herculis&$+1{,}7$&$+0{,}7$\\
13&39 Bootis med.&$+4{,}8$&$+2{,}7$&42&P.\ 17.120&—&$-2{,}1$\\
14&P.\ 14.235&—&$-5{,}0$&43&$\beta$ Draconis&$-0{,}5$&$-1{,}5$\\
15&44 Bootis med.&$+1{,}9$&$+1{,}1$&&&&
\end{tabular}}
\end{center}
\vspace{14mm}
\begin{center}\rule{.25\textwidth}{.4pt}\end{center}
\noindent\hfill $7^*$\par
